\section{Contribución individual de Brenda Lucía Mayhuire Chacón}
\label{sec:brenda-refactorizacion}
\subsection{Responsabilidad y contexto histórico}
Según la identificación y contribución declaradas por Brenda, su responsabilidad es Product Owner y desarrolladora frontend. Como Product Owner relaciona necesidades, prioridades e historias con sus criterios de aceptación; como desarrolladora frontend, su contribución aquí documentada corresponde a HU-01 (registro de usuarios) y HU-02 (asignación de roles y gestión de permisos). Esto no acredita entrevistas, ceremonias ni aceptación de Paz Sur. La distribución inicialmente propuesta en los documentos históricos no se convierte en un acuerdo previo.

La situación inicial descrita al comienzo del capítulo corresponde al momento anterior al primer incremento, sin implementación ejecutable. Posteriormente se creó la versión funcional \texttt{eb00464}, base real de las refactorizaciones. La base pertenece al proyecto y no se atribuye toda V1 a Brenda. No se reconstruye código anterior inexistente.

Se reutilizaron las comparaciones y evidencias de \path{docs/exposicion/brenda/}, conservando su presentación y documentos. Los listados siguientes son extractos literales de \path{src/Users.tsx}, con numeración original y ajuste visual de líneas mediante el estilo \texttt{codigo} del informe. Se incorporan aquí para que no dependan de archivos externos al informe al compilar.

\subsection{HU-01: registro de usuarios, ANTES y DESPUÉS}
\label{sec:brenda-hu01}
En la base, \texttt{submit} preparaba los datos y enviaba la petición. Convertía el documento vacío en \texttt{null}, omitía la contraseña vacía al editar y asignaba roles iniciales solo al registrar. El problema observado era concentrar preparación, transporte y coordinación visual en el mismo bloque; no se demuestra un defecto funcional de ese registro.
\begin{lstlisting}[style=codigo,language=JavaScript,firstnumber=21,caption={HU-01 ANTES: preparación y envío en submit (eb00464).}]
    const payload: Record<string, unknown> = { ...form, document: form.document.trim() || null };
    if (user && !form.password) delete payload.password;
    if (!user) payload.roles = canAssignRoles ? selectedRoles : ['participante'];
    try { await request(`/users${user ? `/${user.id}` : ''}`, { method: user ? 'PUT' : 'POST', body: JSON.stringify(payload) }); onSaved(); }
\end{lstlisting}
En HU-01 se extraen \texttt{prepareUserPayload} y \texttt{saveUser}. El manejador coordina esas llamadas y confirma el éxito solo tras completar la petición; conserva el manejo de errores y el estado de guardado.
\begin{lstlisting}[style=codigo,language=JavaScript,firstnumber=86,caption={HU-01 DESPUÉS: coordinación del guardado (4cab7eb).}]
      const payload = prepareUserPayload(form, user, selectedRoles, canAssignRoles);
      await saveUser(user, payload);
      onSaved();
\end{lstlisting}
La preparación copia el formulario sin modificarlo y conserva las reglas anteriores: omisión de contraseña vacía únicamente al editar y roles iniciales únicamente al crear. La autorización definitiva continua en el servidor.
\begin{lstlisting}[style=codigo,language=JavaScript,firstnumber=34,caption={HU-01 DESPUÉS: preparación del cuerpo de la solicitud (4cab7eb).}]
function prepareUserPayload(
  form: UserFormData,
  user: User | null,
  selectedRoles: string[],
  canAssignRoles: boolean,
): Record<string, unknown> {
  const payload: Record<string, unknown> = {
    ...form,
    document: form.document.trim() || null,
  };

  // En una edicion, no cambiar la contrasena cuando se deja vacia.
  if (user && !form.password) {
    delete payload.password;
  }

  // La asignacion inicial de roles solo corresponde al registro.
  if (!user) {
    payload.roles = canAssignRoles ? selectedRoles : ['participante'];
  }

  return payload;
}
\end{lstlisting}
\begin{lstlisting}[style=codigo,language=JavaScript,firstnumber=59,caption={HU-01 DESPUÉS: transporte HTTP del usuario (4cab7eb).}]
async function saveUser(user: User | null, payload: Record<string, unknown>): Promise<void> {
  await request(`/users${user ? `/${user.id}` : ''}`, {
    method: user ? 'PUT' : 'POST',
    body: JSON.stringify(payload),
  });
}
\end{lstlisting}
Además, \texttt{UserFormData} explicita los campos y \texttt{initialUserForm} reúne los valores iniciales para alta y edición. El estado utiliza inicialización diferida. Esto organiza responsabilidades; no demuestra una mejora de rendimiento.

\subsection{HU-02: asignación de roles y permisos, ANTES y DESPUÉS}
\label{sec:brenda-hu02}
Se muestra el original \texttt{eb00464} para mantener la base común solicitada. El antecedente inmediato de HU-02 es \texttt{4cab7eb}, que conserva las mismas instrucciones de envío de roles y permisos. El diff incremental permite aislar el cambio de HU-02.

Antes, ambos formularios conocían las URLs, el método y la serialización. Después delegan el transporte en auxiliares; conservan la validación de selección obligatoria de roles, los errores y el estado de guardado.
\begin{lstlisting}[style=codigo,language=JavaScript,firstnumber=50,caption={HU-02 ANTES: envío de roles (eb00464).}]
    try { await request(`/users/${user.id}/roles`, { method: 'PUT', body: JSON.stringify({ roles: selected }) }); onSaved(); }
\end{lstlisting}
\begin{lstlisting}[style=codigo,language=JavaScript,firstnumber=105,caption={HU-02 ANTES: envío de permisos (eb00464).}]
    try { await request(`/roles/${role.id}`, { method: 'PUT', body: JSON.stringify({ name, permissions: selected }) }); onSaved(); }
\end{lstlisting}
\begin{lstlisting}[style=codigo,language=JavaScript,firstnumber=174,caption={HU-02 DESPUÉS: delegación de asignación de roles (6ccf4ca).}]
      await saveUserRoles(user.id, selected);
      onSaved();
\end{lstlisting}
\begin{lstlisting}[style=codigo,language=JavaScript,firstnumber=75,caption={HU-02 DESPUÉS: transporte de roles y permisos (6ccf4ca).}]
async function saveUserRoles(userId: User['id'], roles: string[]): Promise<void> {
  await request(`/users/${userId}/roles`, {
    method: 'PUT',
    body: JSON.stringify({ roles }),
  });
}

async function saveRolePermissions(roleId: Role['id'], name: string, permissions: string[]): Promise<void> {
  await request(`/roles/${roleId}`, {
    method: 'PUT',
    body: JSON.stringify({ name, permissions }),
  });
}
\end{lstlisting}
\texttt{RoleForm.submit} utiliza ahora \texttt{saveRolePermissions} con el nombre preparado y los permisos seleccionados. También se extrae \texttt{loadRoleCatalog}, que devuelve ambos catálogos; el componente mantiene la actualización de estado y la protección frente a respuestas tras desmontarse. Las consultas ya usaban \texttt{Promise.all}: no es una optimización nueva.

La selección compartida se concentra en \texttt{toggleSelection}, consumida por los selectores de roles y permisos. Evitar agregar un código ya presente es un ajuste de comportamiento además de reutilización.
\begin{lstlisting}[style=codigo,language=JavaScript,firstnumber=67,caption={HU-02 DESPUÉS: selección reutilizable (6ccf4ca).}]
function toggleSelection(selected: string[], code: string, checked: boolean): string[] {
  if (checked) {
    return selected.includes(code) ? selected : [...selected, code];
  }
  return selected.filter((item) => item !== code);
}
\end{lstlisting}
El nombre del rol se recorta con \texttt{trim}; se rechaza localmente si solo contiene espacios. Antes se enviaba \texttt{name} directamente y ahora se envía \texttt{cleanName}. Este cambio funcional se distingue de la extracción estructural y no reemplaza la validación del servidor.
\begin{lstlisting}[style=codigo,language=JavaScript,firstnumber=263,caption={HU-02 DESPUÉS: nombre validado y guardado de permisos (6ccf4ca).}]
    const cleanName = name.trim();
    if (!cleanName) {
      setError('Escribe un nombre para el rol.');
      return;
    }

    setSaving(true);
    try {
      await saveRolePermissions(role.id, cleanName, selected);
      onSaved();
    } catch (failure) {
\end{lstlisting}
\subsection{Principios aplicados y límites}
\begin{description}
\item[Responsabilidad única.] Preparar datos, transportarlos y coordinar la interfaz tienen funciones identificables. El cuerpo de usuario se localiza en \texttt{prepareUserPayload}; el endpoint, en su función de guardado.
\item[Cohesión.] Cada auxiliar reúne instrucciones de un propósito. El manejador conserva evento, mensajes, guardado y confirmación visual.
\item[Bajo acoplamiento.] Los manejadores dejan de contener URLs y serialización; usan auxiliares con argumentos explícitos. Estos siguen dependiendo de \texttt{request} y del contrato HTTP: el acoplamiento con la API no desaparece.
\item[Reutilización.] \texttt{toggleSelection} tiene consumidores reales en ambos selectores. Las funciones de inicialización, preparación y guardado organizan el formulario compartido por alta y edición.
\end{description}
Las funciones permanecen en \path{src/Users.tsx}; no son módulos nuevos. Los dos commits de refactorización no añaden librerías ni modifican el backend. No se acredita una mejora de velocidad.

Observer podría evaluarse para avisos con consumidores reales y eventos posteriores al commit de la transacción. State sería candidato si los estados de capacitación requieren encapsular comportamientos distintos; usar \texttt{useState} no demuestra ese patrón. Strategy podría servir ante políticas alternativas de evaluación o certificación confirmadas. Ninguno de estos patrones fue implementado por Brenda en los cambios examinados. Extraer funciones auxiliares no equivale a introducir patrones GoF.

\subsection{Pruebas funcionales y alcance de la comprobación}
Brenda reportó comprobaciones manuales de registro de usuario, asignación de roles y edición de permisos, recogidas en \path{docs/exposicion/brenda/GUION.md} y \path{docs/exposicion/brenda/VERIFICACION.md}. Son resultados declarados, sin log nuevo adjunto; no se repitieron esas pruebas manuales en esta adaptación.

El 15 de noviembre de 2026 (America/Lima) se ejecutó nuevamente \texttt{npm.cmd run check}, equivalente aquí a \texttt{npm run check}: \textbf{28 pruebas automáticas aprobadas, 0 fallidas y compilación exitosa} con TypeScript y Vite. Son pruebas del proyecto completo de API, seguridad y persistencia, no 28 pruebas exclusivas del frontend. Usan datos ficticios, SQLite en memoria y bases temporales aisladas, sin utilizar la base de trabajo. Vite advirtió de un paquete mayor de 500 kB sin impedir el build. Véase el registro añadido en \path{docs/PRUEBAS.md}.

Para HU-01 se verifican alta persistida, correo normalizado, duplicados y entradas inválidas. Para HU-02 se comprueban múltiples roles, protección del último administrador, permisos en sesiones existentes y rechazo de escalamiento. Estos casos protegen el contrato del servidor; no ejecutan directamente los auxiliares React ni acreditan todos los flujos del navegador. Esta sesión no ejecuta \texttt{test:ui}. La selección repetida y el nombre con espacios necesitan comprobación específica de interfaz. No se afirma una ejecución previa a cada refactorización ni TDD a partir del resultado actual.

\subsection{Evidencias Git y trazabilidad}
\begin{longtable}{P{2cm}P{2cm}P{8cm}}
\caption{Referencias reales de la contribución documentada de Brenda.}\\
\toprule Commit & Alcance & Evidencia \\ \midrule
\endfirsthead
\toprule Commit & Alcance & Evidencia \\ \midrule
\endhead
\texttt{eb00464} & Base & Primer incremento funcional; original de usuarios, roles y permisos. \\
\texttt{4cab7eb} & HU-01 & Tipado, inicialización, preparación y guardado. Solo modifica \path{src/Users.tsx}. \\
\texttt{6ccf4ca} & HU-02 & Transporte, catálogo, selección compartida, validación y formato JSX. Solo modifica \path{src/Users.tsx}. \\
\bottomrule
\end{longtable}
Comandos de consulta desde la raíz, sin modificar referencias ni código:
\begin{lstlisting}[language={},numbers=none,caption={Consultas reproducibles de la evidencia Git.}]
git show eb00464:src/Users.tsx
git diff eb00464 4cab7eb -- src/Users.tsx
git diff 4cab7eb 6ccf4ca -- src/Users.tsx
git diff eb00464 6ccf4ca -- src/Users.tsx
\end{lstlisting}
Los doce fragmentos preservados en la exposición se cotejaron con Git, sus rangos y las huellas SHA-256 de \path{docs/exposicion/brenda/MANIFIESTO.json}, sin modificarlos. Los commits de refactorización tienen un nombre de autor genérico: la identificación de la contribución de Brenda procede de su declaración y del encargo, no de ese nombre. No se alteran identidades Git.

La revisión técnica asistida por IA comprende lectura del diff, cotejo de extractos y comprobaciones registradas. No equivale a revisión humana por otro integrante, aceptación de Paz Sur ni integración mediante merge. Conforme a la Definition of Done, las pruebas no bastan para declarar Hechas las historias completas. No se realizaron commits, merge ni push en esta adaptación.
